\chapter{Factorizations of a Joint Probability Distribution}
\label{ch-factor-joint-prob}

This chapter is based
entirely on Question/Answer sessions with ChatGPT, reported in my own words and notation.
Also, I added a small amount of personal opinions and vision.

In this chapter, we discuss 3 important
ways of factoring a joint probability distribution:
\begin{enumerate}
\item Factorization into families
\item Factorization into max-cliques
\item FED factorization. (FED=Family
of exponential Distributions)

\end{enumerate}
We analyze the same bnet (the on in Fig.\ref{fig-2-boxes-bnet})
in those 3 factorizations. We give
various pictorials representations (PR) of those factorizations.

\section{Factorization into Families}

We will discuss 2 types of factorizations
into families,
Bayesian networks (bnets),
and family hypergraphs.




\begin{itemize}
\item {\bf Bayesian Network}

Factorization  of bnet in Fig.\ref{fig-2-boxes-bnet}:
\beq
P(x^6)
=
P(a)
\quad P(b|a)
\quad P(c|a)
\quad P(d|b,c)
\quad P(e|d)
\quad P(f|c,e)
\eeq

\begin{figure}[h!]
$$
\begin{array}{ccc}
\xymatrix{
\rva\ar[r]\ar[d]
&{\rvb}\ar[d]
\\
{\rvc}\ar[r]\ar[d]
&{\rvd}\ar[d]
\\
\rvf
&\rve\ar[l]
}
&&
\xymatrix{
\rva
\ar@{-}[r]
\ar@{-}[d]
&{\rvb}
\ar@{-}[d]
\ar@{-}[dl]
\\
{\rvc}
\ar@{-}[r]
\ar@{-}[d]
\ar@{-}[dr]
&{\rvd}
\ar@{-}[d]
\\
\rvf
&\rve
\ar@{-}[l]
}
\\
(a)
&& (b)
\end{array}
$$
\caption{Example of a bnet $(a)$ and
its moral graph $(b)$. Let
$x^6 = (a,b,c,d,e, f)$.
Also, we will refer to bnet $(a)$ as $\calb$ and to the set
$\{\rva, \rvb, \rvc, \rvd, \rve, \rvf\}$ of $\calb$'s nodes as $nds(\calb)$.}
\label{fig-2-boxes-bnet}
\end{figure}

PR of this factorization in Fig.\ref{fig-2-boxes-bnet}.

\item {\bf Family Hypergraph}

Let
\beq
fam(\rvx)= \{\rvx\}
\cup pa(\rvx)
\eeq

Factorization  of bnet in Fig.\ref{fig-2-boxes-bnet}:


\begin{align}
P(x^6)&=
\frac{1}{Z}
\prod_{\rvx\in nds(\calb)}
\phi(fam(\rvx))
\\
&=
\frac{1}{Z}
\phi_1(a)
\quad \phi_2(b,a)
\quad \phi_3(c,a)
\quad \phi_4(d, b, c)
\quad \phi_5(e, d)
\quad \phi_6(f, c, e)
\end{align}
The Ising model is an example of
factorization into family hypergraph.

PR of this factorization in Fig.\ref{fig-tab-fam-hyper}, Fig.\ref{fig-fact-g-fam-hyper}, and
Fig.\ref{fig-facet-fam-hyper}


\begin{figure}[h!]
\begin{tabular}{|l|l|l|l|l|l|l|}
\hline
 & \cellcolor[HTML]{FFCCC9}$\phi_1$ & \cellcolor[HTML]{FFCCC9}$\phi_2$ & \cellcolor[HTML]{FFCCC9}$\phi_3$ & \cellcolor[HTML]{FFCCC9}$\phi_4$ & \cellcolor[HTML]{FFCCC9}$\phi_5$ & \cellcolor[HTML]{FFCCC9}$\phi_6$ \\ \hline
\cellcolor[HTML]{9AFF99}$\rva$ & $\color{red}\bullet$ & $\color{red}\bullet$ & $\color{red}\bullet$ &  &  &  \\ \hline
\cellcolor[HTML]{9AFF99}$\rvb$ &  & $\color{red}\bullet$ &  & $\color{red}\bullet$ &  &  \\ \hline
\cellcolor[HTML]{9AFF99}$\rvc$ &  &  & $\color{red}\bullet$ & $\color{red}\bullet$ &  & $\color{red}\bullet$ \\ \hline
\cellcolor[HTML]{9AFF99}$\rvd$ &  &  &  & $\color{red}\bullet$ & $\color{red}\bullet$ &  \\ \hline
\cellcolor[HTML]{9AFF99}$\rve$ &  &  &  &  & $\color{red}\bullet$ & $\color{red}\bullet$ \\ \hline
\cellcolor[HTML]{9AFF99}$\rvf$ &  &  &  &  &  & $\color{red}\bullet$ \\ \hline
\end{tabular}
\caption{Table PR of family hypergraph.}
\label{fig-tab-fam-hyper}
\end{figure}





\begin{figure}[h!]
$$
\xymatrix{
\phi_1\ar[d]
&\phi_2\ar[d]\ar[dl]
&\phi_3\ar[d]\ar[dll]
&\phi_4\ar[dll]\ar[dl]\ar[d]
&\phi_5\ar[d]\ar[dl]
&\phi_6\ar[dl]\ar[d]\ar[dll]
\\
\rva
&\rvb
&\rvc
&\rvd
&\rve
&\rvf
}
$$
\caption{Factor graph (bipartite) PR of
family hypergraph.}
\label{fig-fact-g-fam-hyper}
\end{figure}

\begin{figure}
$$
\xymatrix@C=3pc@R=3pc{
\phi_1,\rva
\ar@{-}[r]|{\phi_2}
\ar@{-}[d]|{\phi_3}
&{\rvb}
\ar@{-}[d]|{\phi_4}
\ar@{-}[dl]|{\phi_4}
\\
{\rvc}
\ar@{-}[r]|{\phi_4}
\ar@{-}[d]|{\phi_6}
\ar@{-}[dr]|{\phi_6}
&{\rvd}
\ar@{-}[d]|{\phi_5}
\\
\rvf
&\rve
\ar@{-}[l]|{\phi_6}
}
$$
\caption{Facet PR of family hypergraph.}
\label{fig-facet-fam-hyper}
\end{figure}

\end{itemize}

\section{Factorization into Max-Cliques}

We will
only discuss only one type of factorization
into max-cliques, {Markov Random Fields (MRF)}.

A {\bf clique} of a DAG (directed acyclic graph)
or an UG (undirected graph) is a fully connected subgraph. A {\bf maximal clique (max-clique)} is a clique which is not contained in a larger clique.

To construct a {\bf Markov Random Field (MRF)} from a given bnet, one performs the
following steps

\begin{enumerate}
\item Find the moral graph of the bnet (see Chapter \ref{ch-dsep-ug}).
\item Find the max-cliques of the moral graph.
\item Write down a factorization into
max-cliques.
\end{enumerate}


Fig.\ref{fig-2-boxes-bnet} gives an example of a bnet and its moral graph.
The max-cliques for that moral graph are

\beq
\rva\rvb\rvc, \rvb\rvc\rvd, \rvc\rvd\rve, \rvc\rve\rvf
\eeq
where, for example, $\rva\rvb\rvc$
is an abbreviation for $\{\rva, \rvb, \rvc\}$.

Factorization  of bnet in Fig.\ref{fig-2-boxes-bnet}:

\beqa
P(x^6)&=&
\frac{1}{Z}
\prod_{C\in cliques(\calb)}
\phi(C)
\\
&=&
\frac{1}{Z}
\phi_1(a,b,c)
\quad\phi_2(b,c,d)
\quad\phi_3(c,d,e)
\quad\phi_4(c,e,f)
\eeqa

PR of this factorization in Fig.\ref{fig-tab-mrf},
Fig.\ref{fig-fact-g-mrf}, and
Fig.\ref{fig-facet-mrf}

\begin{figure}[h!]
\centering
\begin{tabular}{lllll}
\hline
\multicolumn{1}{|l|}{} & \multicolumn{1}{l|}{\cellcolor[HTML]{FFCCC9}$\phi_1$} & \multicolumn{1}{l|}{\cellcolor[HTML]{FFCCC9}$\phi_2$} & \multicolumn{1}{l|}{\cellcolor[HTML]{FFCCC9}$\phi_3$} & \multicolumn{1}{l|}{\cellcolor[HTML]{FFCCC9}$\phi_4$} \\ \hline
\multicolumn{1}{|l|}{\cellcolor[HTML]{9AFF99}$\rva$} & \multicolumn{1}{l|}{$\color{red}\bullet$} & \multicolumn{1}{l|}{} & \multicolumn{1}{l|}{} & \multicolumn{1}{l|}{} \\ \hline
\multicolumn{1}{|l|}{\cellcolor[HTML]{9AFF99}$\rvb$} & \multicolumn{1}{l|}{$\color{red}\bullet$} & \multicolumn{1}{l|}{$\color{red}\bullet$} & \multicolumn{1}{l|}{} & \multicolumn{1}{l|}{} \\ \hline
\multicolumn{1}{|l|}{\cellcolor[HTML]{9AFF99}$\rvc$} & \multicolumn{1}{l|}{$\color{red}\bullet$} & \multicolumn{1}{l|}{$\color{red}\bullet$} & \multicolumn{1}{l|}{$\color{red}\bullet$} & \multicolumn{1}{l|}{$\color{red}\bullet$} \\ \hline
\multicolumn{1}{|l|}{\cellcolor[HTML]{9AFF99}$\rvd$} & \multicolumn{1}{l|}{} & \multicolumn{1}{l|}{$\color{red}\bullet$} & \multicolumn{1}{l|}{$\color{red}\bullet$} & \multicolumn{1}{l|}{} \\ \hline
\multicolumn{1}{|l|}{\cellcolor[HTML]{9AFF99}$\rve$} & \multicolumn{1}{l|}{} & \multicolumn{1}{l|}{} & \multicolumn{1}{l|}{$\color{red}\bullet$} & \multicolumn{1}{l|}{$\color{red}\bullet$} \\ \hline
\multicolumn{1}{|l|}{\cellcolor[HTML]{9AFF99}$\rvf$} & \multicolumn{1}{l|}{} & \multicolumn{1}{l|}{} & \multicolumn{1}{l|}{} & \multicolumn{1}{l|}{$\color{red}\bullet$} \\ \hline
 &  &  &  &
\end{tabular}
\caption{Table PR of MRF}
\label{fig-tab-mrf}
\end{figure}


\begin{figure}[h!]
$$
\xymatrix{
\phi_1
\ar[d]\ar[dr]\ar[drr]
&\phi_2
\ar[d]\ar[dr]\ar[drr]
&\phi_3\ar[d]\ar[dr]\ar[drr]
&\phi_4\ar[dr]\ar[drr]\ar[dl]
\\
\rva
&\rvb
&\rvc
&\rvd
&\rve
&\rvf
}
$$
\caption{Factor graph (bipartite)
PR of MRF}
\label{fig-fact-g-mrf}
\end{figure}

\begin{figure}[h!]
$$
\xymatrix@C=4pc@R=4pc{
\rva
\ar@{-}[r]|{\phi_1}
\ar@{-}[d]|{\phi_1}
&{\rvb}
\ar@{-}[d]|{\phi_2}
\ar@{-}[dl]|{\phi_1}
\ar@{-}@/_1pc/[dl]|{\phi_2}
\\
{\rvc}
\ar@{-}[r]|{\phi_2}
\ar@/_1pc/@{-}[r]|{\phi_3}
\ar@{-}[d]|{\phi_4}
\ar@{-}[dr]|{\phi_3}
\ar@{-}@/_1pc/[dr]|{\phi_4}
&{\rvd}
\ar@{-}[d]|{\phi_3}
\\
\rvf
&\rve
\ar@{-}[l]|{\phi_4}
}
$$
\caption{Facet PR of MRF }
\label{fig-facet-mrf}
\end{figure}
\section{FED Factorization}

Let

FED= Family of Exponential Distributions.\footnote{FEDs are discussed and used in Section \ref{sec-fed}.}

Factorization  of bnet in Fig.\ref{fig-2-boxes-bnet}:

\beq
P(x^6)=\frac{1}{Z}
\prod_{i=1}^{15}
\exp\left(\cals_i(x^6)\theta_i
\right)
\eeq
so

\beq
\ln P(x^6)=
\sum_{i=1}^{15} \cals_i(x^6)\theta_i
-\ln Z
\label{eq-fed-fact-ln}
\eeq
were the {\bf statistics} $\cals_i(x^6)$ will be specified below.
Eq.(\ref{eq-fed-fact-ln})
is analogous to an
 expansion of $\ln P(x^6)$
in terms  of a complete or over-complete orthonormal set, such a Fourier series. The
complete set is finite because
$\ln P(x^6)$ is finite and discrete and thus has a finite number of degrees of freedom.

PR of this factorization in Fig.\ref{fig-bnet-fed} and
Fig.\ref{fig-facet-fed}.
\begin{figure}[h!]
$$
\xymatrix@C=4pc@R=1pc{
&\ul{\theta_1}\ar[ddl]|{\cals_1(x^6)}
\\
&\ul{\theta_2}\ar[dl]|{\cals_2(x^6)}
\\
\ln P(x^6)&\vdots
\\
&\ul{\theta_{15}}\ar[ul]|{\cals_{14}(x^6)}
\\
&\ul{\theta_{15}}\ar[uul]|{\cals_{15}(x^6)}
}$$
\caption{bnet PR of the FED factorization.}
\label{fig-bnet-fed}
\end{figure}

\begin{figure}[h!]
$$
\begin{array}{ccc}
\xymatrix@C=3pc@R=3pc{ % nodes
\rva, \phi_1
\ar@{-}[r]
\ar@{-}[d]
&{\rvb},\phi_2
\ar@{-}[d]
\ar@{-}[dl]
\\
{\rvc},\phi_3
\ar@{-}[r]
\ar@{-}[d]
\ar@{-}[dr]
&{\rvd},\phi_4
\ar@{-}[d]
\\
\rvf,\phi_6
&\rve,\phi_5
\ar@{-}[l]
}
&
\xymatrix@C=3pc@R=3pc{% duos
\rva
\ar@{-}[r]|{\phi_1'}
\ar@{-}[d]|{\phi_2'}
&{\rvb}
\ar@{-}[d]|{\phi_4'}
\ar@{-}[dl]|{\phi_3'}
\\
{\rvc}
\ar@{-}[r]|{\phi_5'}
\ar@{-}[d]|{\phi_6'}
\ar@{-}[dr]|{\phi_7'}
&{\rvd}
\ar@{-}[d]|{\phi_8'}
\\
\rvf
&\rve
\ar@{-}[l]|{\phi_9'}
}
&
\xymatrix@C=3pc@R=3pc{% trios
\rva
\ar@{-}[r]
\ar@{-}[d]
&{\rvb}
\ar@{-}[d]|{\phi_1''}
\ar@{-}[dl]|{\phi_1''}
\\
{\rvc}
\ar@{-}[r]|{\phi_1''}
\ar@{-}@/_1pc/[r]|{\phi_2''}
\ar@{-}[d]
\ar@{-}[dr]|{\phi_2''}
&{\rvd}
\ar@{-}[d]|{\phi_2''}
\\
\rvf
&\rve
\ar@{-}[l]
}
\\
(a)&(b)&(c)
\end{array}
$$
\caption{Facet PR of FED
factorization.
Normally, we would super-pose $(a), (b)$ and $(c)$, but for increased clarity, we
have opted to draw single, double,
and triple interactions separately.
Each $\phi_i$ factor equals $\exp(\cals_i(x^6)\theta_i)$.
$(a)$ Self-interactions (6 of them),
$(b)$ pair interactions (9 of them),
(3) triple interactions (2 of them). Total
number of interactions shown: $6+9+2=17$.}
\label{fig-facet-fed}
\end{figure}
Let

TPM = Transition probability matrix

 dofs=degrees of freedom

The family of statistics functions  $\{\cals_i(x^6)\}_{i=1}^{17}$
 is
far from unique. We will give two natural choices:
TPM dofs statistics and interaction statistics.

\hrule
{\bf TPM dofs statistics:}

\beq
a=0:\quad \cals_1(x^6)=\delta^{a}_0,
\eeq

\beq
b=0, \text{vary } a:\quad \cals_2(x^6)=\delta^{b}_0 \delta^a_0,
\quad
\cals_3(x^6)=\delta^{b}_0 \delta^a_1
\eeq

\beq
c=0, \text{vary } a:\quad
\cals_4(x^6)=\delta^{c}_0 \delta^a_0,
\quad
\cals_5(x^6)=\delta^{c}_0 \delta^a_1
\eeq

\beq
\begin{array}{l}
d=0, \text{vary } b, c:\quad
\\
\cals_6(x^6)=\delta^{d}_0
\delta_0^b \delta_0^c,
\quad
\cals_7(x^6)=\delta^{d}_0
\delta_0^b \delta_1^c,
\quad
\cals_8(x^6)=\delta^{d}_0
\delta_1^b \delta_0^c,
\quad
\cals_9(x^6)=\delta^{d}_0
\delta_1^b \delta_1^c
\end{array}
\eeq

\beq
e=0, \text{vary } d:\quad
\cals_{10}(x^6)=\delta^{e}_0 \delta^d_0,
\quad
\cals_{11}(x^6)=\delta^{e}_0 \delta^d_1
\eeq

\beq
\begin{array}{l}
f=0, \text{vary } c, e:\quad
\\
\cals_{12}(x^6)=\delta^{f}_0
\delta_0^c \delta_0^e,
\quad
\cals_{13}(x^6)=\delta^{f}_0
\delta_0^c \delta_1^e
,\quad
\cals_{14}(x^6)=\delta^{f}_0
\delta_1^c \delta_0^e,
\quad
\cals_{15}(x^6)=\delta^{f}_0
\delta_1^c \delta_1^e
\end{array}
\eeq

Recall that
\beq
P(d, b|c) = P(d|b, c)P(b|c)
\eeq
We've parametrized $P(d|b,c)$ so far, but
we also want to parametrize $P(b|c)$.
So let's include the additional statistic

\beq
\cals_{16}(x^6)= \delta_0^b\delta_0^c
\eeq

Likewise, recall that
\beq
P(f,c|e) = P(f|c,e)P(c|e)
\eeq
We've parametrized $P(f|c,e)$ so far, but
we also want to parametrize $P(c|e)$.
So let's include the additional statistic

\beq
\cals_{17}(x^6)= \delta_0^c\delta_0^e
\eeq
Note that the statistics $\cals_{16}$
and $\cals_{17}$ can also
be understood as coming from the two new
edges that arise from moralization.

\hrule
{\bf Interaction statistics}

Consider the set of interactions
\beq
\cali=
\begin{array}{l}
\{
\eset, \rva, \rvb, \rvc, \rvd, \rve, \rvf,
\\
\rva\rvb, \rva\rvc,
\rvb\rvc,\rvb\rvd,
\rvc\rvd, \rvd\rve,\rvc\rve,
\rvc\rvf, \rve\rvf,
\\
\rvb\rvc\rvd, \rvc\rve\rvf
\}
\end{array}
\eeq
where, for example, $\rva\rvb\rvc$
is an abbreviation for $\{\rva, \rvb, \rvc\}$.
Now the FED
factorization equation Eq.(\ref{eq-fed-fact-ln})
becomes

\beq
\ln P(x^6)=
\sum_{A \in \cali}
 \cals_A(x^6)\theta_A
\eeq
where we've set

\beq
\theta_\eset =
-\ln Z,
\quad
\cals_\eset(x^6)=1
\eeq

Define
\beq
0^A 1^{A^c}=
(\underbrace{0\ldots 0}_{A},
\underbrace{1,\dots, 1}_{A^c})
\eeq
For example,

\beq
0^{\rvb\rvc\rvd}1^{\rvb\rvc\rvd^c}
=
\begin{tabular}{cccccc}
$\rva$&$\rvb$&$\rvc$&$\rvd$&$\rve$&$\rvf$
\\
(1,&0,&0,&0,&1,&1)
\end{tabular}
\eeq
Also define

\beq
\ell_A = \ln P(0^A 1^{A^c})
\eeq
Now the FED factorization
equation Eq.(\ref{eq-fed-fact-ln}) can be
expressed compactly
as the matrix equation

\beq
\ell = M \theta
\eeq
where $M$ is an $18\times 18$  matrix.

To calculate what are the
matrix elements of matrix $M$, note the following.
For $A\subset nds(\calb)$,

\beq
\cals_A(x^6)=\prod_{\rvx\in A}\delta_0^x=
\delta_0^{x_A}
\eeq
Thus

\beq
\ln P(0^B 1^{B^c})=\sum_{A\in \cali}\delta_0^{x_A}\indi
(x=0^B 1^{B^c})\theta_A=
\sum_{A\in \cali}
\underbrace{\indi(A\subset B)}_
{M_{B,A}}
\theta_A
\eeq
Thus

\beq
M_{B,A}= \indi(A\subset B)
\eeq
Note that $M$ is lower triangular
with all elements along its diagonal
equal to 1,
and all elements
under its diagonal equal
to zero or one.

As a consequence of the M\"{o}bius inversion
formula\footnote{The M\"{o}bius inversion
formula is discussed and proven in Section \ref{sec-mobius-inv}.}

\beq
\theta_B =
\sum_{A\in \calp(nds(\calb))}
\indi(A \subset B)
(-1)^{|A|-|B|}\ell_A
\label{eq-mob-inv-interaction-case}
\eeq
ChatGPT evaluated Eq.(\ref{eq-mob-inv-interaction-case}) using the Python package
for symbolic algebra,
SymPy, and assures me that the solution is as follows.
For $\rvx, \rvy, \rvz\in nds(\calb)$,

{\bf self interactions (singletons)}
\beq
\theta_\rvx=\ell_\rvx-\ell_\eset
\eeq

{\bf pair interactions}
\beq
\theta_{\rvx\rvy}=\ell_{\rvx\rvy}
-(\ell_\rvx
+\ell_\rvy)
+\ell_\eset
\eeq

{\bf triple interactions}

\beq
\theta_{\rvx\rvy\rvz}=
\ell_{\rvx\rvy\rvz}-(\ell_{\rvx\rvy}
+\ell_{\rvy\rvz}
+\ell_{\rvx\rvz}
)+(\ell_\rvx + \ell_\rvy +\ell_\rvz)
-\ell_\eset
\eeq

\begin{claim}
\label{cl-geo-series}
If $M = I + N\in \RR^{n\times n}$ is a lower triangular matrix so that $N$ is strictly lower triangular, then

\beq
N^n=0\quad \text{($N$ is {\bf nilpotent})}
\eeq
and
\beq
M^{-1}=I -N +N^2 -\ldots N^n
\eeq
\end{claim}
\proof

$N^n=0$ is obvious. Just work out an example with $n=3$ to convince yourself.
It can be proven by induction if you are a masochist.

As for the formula for $M^{-1}$, suppose $n=4$. Then
\beq
(1+N)(1-N+N^2-N^3)=
1-N^4=1
\eeq
This argument is clearly valid for $n\geq 2$.
\qed

We won't do it here, but note that one can use Claim
\ref{cl-geo-series}
instead of the M\"{o}bius inversion
formula (MIF)
to invert the $M$
 matrix in $\ell=M\theta$.
In fact, Claim
\ref{cl-geo-series} is more general
than the MIF and can be used for other statistic families such as our TPM dofs statistics
for which MIF does not directly apply.

Note that
the bnet of Eq.(\ref{fig-2-boxes-bnet})  has 15 dofs
but the column vectors $\ell$ and $\theta$
have length 18. So our interaction statistics
for this bnet yields
an over complete expansion of $
\ln P(x^6)$.
One of the extra 3 statistics is
being used to accommodate $\eset$ and $\ln Z$. The other two extra
statistics come from
moralizing the parents of colliders. Thus, the bnet of
Fig.\ref{fig-2-boxes-bnet}
is represented by a 15 dim submanifold of the larger
18 dim manifold. The submanifold satisfies the
following two conditional
independence constraints.

\beq
\rvb \perp\rvc \mid \rvd,
\quad
\rvc\perp \rve \mid \rvf
\eeq


\begin{claim}
Suppose $\rvx, \rvy, \rvz$ are binary random variables. If
$\rvx\perp \rvz\mid \rvy$ then

\beq
\frac{P_{\rvx,\rvy,\rvz}(0,y,0)}
{P_{\rvx,\rvy,\rvz}(1,y,0)}
=
\frac{P_{\rvx,\rvy,\rvz}(0,y,1)}
{P_{\rvx,\rvy,\rvz}(1,y,1)}
\eeq

\end{claim}
\proof

\beq
P_{\rvx|\rvz, \rvy}(x|z, y)=
P_{\rvx|\rvy}(x|y)
\eeq

\beq
P_{\rvx|\rvz, \rvy}(x|0, y)=
P_{\rvx|\rvz, \rvy}(x|1, y)=
P_{\rvx|\rvy}(x|y)
\eeq

\beq
P_{\rvx|\rvz, \rvy}(0|0, y)=
P_{\rvx|\rvz, \rvy}(0|1, y)=
P_{\rvx|\rvy}(0|y)
\eeq

\beq
P_{\rvx|\rvz, \rvy}(1|0, y)=
P_{\rvx|\rvz, \rvy}(1|1, y)=
P_{\rvx|\rvy}(1|y)
\eeq

\beq
\frac{P_{\rvx|\rvz, \rvy}(0|0, y)}
{P_{\rvx|\rvz, \rvy}(1|0, y)}
=
\frac{P_{\rvx|\rvz, \rvy}(0|1, y)}
{P_{\rvx|\rvz, \rvy}(1|1, y)}
\eeq
\qed

Let

$Q=(\rvx,\rvy, \rvz)$,


$Q^c= nds(\calb)-Q$

Then

\beq
P_{\rvx,\rvy, \rvz}(x,y,z)=
\sum_{x_{Q^c}\in val(Q^c)}
\exp\left(  \sum_{A\in \cali}
\cals_A(x^6)\theta_A\right)
\eeq


Two constraints:

\begin{enumerate}

\item $\rvb\perp\rvc\mid\rvd$

\beq
\frac{P_{\rvb,\rvd,\rvc}(0,d,0)}
{P_{\rvb,\rvd,\rvc}(1,d,0)}
=
\frac{P_{\rvb,\rvd,\rvc}(0,d,1)}
{P_{\rvb,\rvd,\rvc}(1,d,1)}
\eeq

\item $\rvc\perp\rve \mid \rvf$


\beq
\frac{P_{\rvc,\rvf,\rve}(0,f,0)}
{P_{\rvc,\rvf,\rve}(1,f,0)}
=
\frac{P_{\rvc,\rvf,\rve}(0,f,1)}
{P_{\rvc,\rvf,\rve}(1,f,1)}
\eeq
\end{enumerate}
\hrule
Simplicial Complexes

Given a set of vertices $V$ with powerset $\calp(V)$,
a {\bf hypergraph}
$\calh$ is a subset
of $\calp(V)$. Thus, $\calh\subset \calp(V)$.

A {\bf simplicial complex} $\calh$
is a hypergraph
that has the {\bf downward-closed (or hereditary)} property:
If $A\in \calh$ and $B\subset A$, then $B\in \calh$.
In particular, $\eset \in \calh$.
The elements of $\calh$ are called {\bf
simplices or simplexes}.
(I've been calling them interactions or facets)


The hypergraph of Fig.\ref{fig-facet-fed} is an example of a simplicial complex if we add to it the empty set.
To obtain the hypergraph $\cali$ represented in that figure,
we started with the families of the bnet, and generated a simplicial complex by
enforcing the hereditary property.
This generated 2 out of the 4 max-cliques
of the moral graph. The two max-cliques
that were generated were the ones
introduced by moralizing colliders.
